\section*{Software Requirements Specification}

\section{Introduction}

\subsection{Purpose}
Настоящий документ описывает функциональные и нефункциональные требования
к разработке веб-портала «Календарь.ру» — информационного ресурса,
предоставляющего пользователям доступ к календарным данным, праздникам,
именинам, гороскопам и историческим событиям.

\subsection{Scope}
Документ распространяется на разработку веб-портала,
включающего модули календаря, праздников, именин, астрологических прогнозов,
исторических событий, редакционного контента, рекламных размещений, поиска
и пользовательской персонализации. Система
предназначена для массовой аудитории пользователей интернета в возрасте
от 16 до 65 лет.

\subsection{Definitions, Acronyms and Abbreviations}
\begin{enumerate}[leftmargin=*,label=\arabic*.]
    \item SRS — Software Requirements Specification, спецификация требований к программному обеспечению
    \item PWA — Progressive Web Application, прогрессивное веб-приложение
    \item RUP — Rational Unified Process, рациональный унифицированный процесс разработки ПО
    \item WCAG — Web Content Accessibility Guidelines, руководящие принципы доступности веб-контента
    \item CDN — Content Delivery Network, сеть доставки контента для ускорения загрузки ресурсов
    \item XSS — Cross-Site Scripting, уязвимость, позволяющая внедрять вредоносный код в веб-страницы
    \item CSRF — Cross-Site Request Forgery, атака подделки межсайтовых запросов
    \item SQL-инъекция — уязвимость, позволяющая внедрять произвольные SQL-запросы в приложение
    \item bcrypt — криптографический алгоритм хэширования паролей с защитой от брутфорса
    \item SEO — Search Engine Optimization, поисковая оптимизация страниц и контента
    \item CTR — Click-Through Rate, показатель кликабельности рекламного блока
\end{enumerate}


\subsection{References}
\begin{itemize}[leftmargin=*]
    \item Официальный сайт: \url{https://www.calend.ru/}
    \item ГОСТ Р ИСО/МЭК/ИЕЕЭ 29148-2019 «Системная и программная инженерия»
\end{itemize}

\subsection{Overview}
Документ состоит из семи основных разделов, последовательно раскрывающих
требования к веб-порталу «Календарь.ру». Разделы 2–5 содержат детальное
описание функциональных и нефункциональных требований, внешних интерфейсов
и ограничений системы. Разделы 6–7 визуализируют и описывают сценарии
взаимодействия пользователей с системой через диаграммы и текстовые
прецеденты. Приложения включают оценку трудоёмкости требований и анализ
проектных рисков для эффективного управления разработкой.

\section{Overall Description}

\subsection{Product Features}
Разрабатываемая система представляет собой информационный веб-портал
"Календарь.ру", предоставляющий пользователям доступ к
календарным данным. Основные особенности продукта:
\begin{itemize}[leftmargin=*]
    \item Интерактивный календарь с возможностью навигации по месяцам и годам (1900–2100 гг.)
    \item Информация о государственных, профессиональных, народных и религиозных праздниках с историческими справками
    \item Календарь именин с указанием значений имён, происхождения и дат празднования по православному и католическому календарям
    \item Астрологические прогнозы: ежедневные, еженедельные и ежемесячные гороскопы для 12 знаков зодиака с разделением по сферам (любовь, финансы, здоровье)
    \item Редакционный раздел с новостями, статьями и анонсами, включая черновики и отложенную публикацию
    \item Поддержка рекламных блоков и партнёрских материалов с настройкой размещения и базовой аналитикой
\end{itemize}

\subsection{User Characteristics}
Система поддерживает четыре категории пользователей с различными правами доступа:
\begin{itemize}[leftmargin=*]
    \item \textbf{Гость} — анонимный посетитель без учётной записи:
    \begin{itemize}
        \item Просмотр календаря, праздников, именин, гороскопов и исторических событий
        \item Поиск контента по ключевым словам, датам и категориям
        \item Просмотр новостей и анонсов на главной странице
    \end{itemize}
    \item \textbf{Зарегистрированный пользователь} — пользователь с подтверждённой учётной записью:
    \begin{itemize}
        \item Все возможности гостя
        \item Добавление контента в избранное и создание персональных закладок
        \item Настройка персонального гороскопа (выбор знака зодиака)
        \item Получение уведомлений о праздниках и именинах по электронной почте
        \item Сохранение истории поиска
    \end{itemize}
    \item \textbf{Редактор контента} — сотрудник редакции с правами на управление информацией:
    \begin{itemize}
        \item Все возможности зарегистрированного пользователя
        \item Создание, редактирование и публикация новостей, статей и тематических подборок
        \item Добавление, редактирование и удаление праздников, именин, гороскопов и исторических событий
        \item Загрузка изображений и медиаконтента
        \item Модерация пользовательских материалов (при наличии функционала комментариев или обратной связи)
        \item Управление рекламными размещениями в доступных разделах портала
        \item Просмотр статистики просмотров материалов и базовой рекламной статистики
    \end{itemize}
    \item \textbf{Администратор системы} — технический специалист с полными правами доступа:
    \begin{itemize}
        \item Все возможности редактора контента
        \item Управление пользователями (блокировка, сброс паролей, изменение ролей)
        \item Настройка параметров системы, интеграций и рекламных платформ
        \item Мониторинг работоспособности и производительности
        \item Управление резервным копированием и восстановлением данных
        \item Анализ полной статистики посещаемости сайта и эффективности рекламных кампаний
    \end{itemize}
\end{itemize}

\subsection{Operating Environment}
\begin{itemize}[leftmargin=*]
    \item \textbf{Клиентская часть:} Веб-браузеры Chrome 100+, Firefox 95+, Safari 14+, Edge 95+; мобильные устройства iOS 14+, Android 10+; адаптивный дизайн от 320px.
    \item \textbf{Серверная часть:} Операционная система — Linux (Ubuntu 22.04 LTS); веб-сервер — Nginx 1.20+; backend — Java 17+ с Spring Boot 3.x; база данных — PostgreSQL 14+; кэширование — Redis 7+; поиск — Elasticsearch 8+.
    \item \textbf{Инфраструктура:} Поддержка HTTPS/TLS 1.3; интеграция с CDN для статики; мониторинг через Prometheus + Grafana.
\end{itemize}

\subsection{Assumptions and Dependencies}
\begin{description}[leftmargin=!,labelwidth=\widthof{\bfseries}]
    \item[2.4.1] Клиентская часть реализована с использованием React, HTML5, CSS3 и JavaScript (ES6+). Веб-сайт должен корректно отображаться в браузерах: Chrome 100+, Firefox 95+, Safari 14+, Edge 95+.
    \item[2.4.2] Серверная часть реализована на Java 17+ с использованием Spring Boot 3.x. В качестве веб-сервера используется Nginx 1.20+. Для хранения данных применяется PostgreSQL 14+ (пользователи и основной контент) и Redis 7+ (кэширование). Серверное ПО развертывается на операционной системе Linux (Ubuntu 22.04 LTS).
\end{description}

\subsection{Constraints}
Обмен данными между клиентом и сервером происходит по протоколу HTTPS с
использованием TLS 1.3, формат данных — JSON. Для отправки почтовых уведомлений 
используется SMTP-сервис. Все пароли пользователей хранятся в
хэшированном виде с использованием алгоритма bcrypt. Персональные данные обрабатываются
и хранятся на серверах, расположенных на
территории Российской Федерации. Для персонализированной рекламы и аналитики
должно использоваться управление пользовательскими согласиями (cookie consent).

\section{System Features}
\subsection{Functional Requirements}
\subsubsection*{Календарные данные и тематический контент}
\begin{description}[leftmargin=!,labelwidth=\widthof{\bfseries}]
    \item[3.1.1] Система должна отображать календарь текущего месяца при загрузке главной страницы.
    \item[3.1.2] Система должна позволять навигацию по месяцам (вперёд/назад).
    \item[3.1.3] Система должна позволять переход к любому году в диапазоне 1900–2100.
    \item[3.1.4] Система должна обозначать выходные дни отдельно от рабочих дней.
    \item[3.1.5] Система должна обозначать праздничные дни отдельным признаком.
    \item[3.1.6] Система должна указывать текущую календарную дату.
    \item[3.1.7] Система должна отображать номер недели в отдельной колонке.
    \item[3.1.8] Система должна отображать список праздников на выбранную дату с указанием типа (государственный, профессиональный, народный, религиозный).
    \item[3.1.9] Система должна отображать подробную информацию о празднике: историю, традиции, статус (выходной/рабочий день).
    \item[3.1.10] Система должна позволять поиск праздников по названию или части названия.
    \item[3.1.11] Система должна позволять фильтрацию праздников по типу и году.
    \item[3.1.12] Система должна отображать календарь праздников на весь год в табличном виде.
    \item[3.1.13] Система должна предоставлять информацию о переносе выходных дней согласно постановлениям правительства.
\end{description}

\subsubsection*{Поиск, авторизация и персонализация}
\begin{description}[leftmargin=!,labelwidth=\widthof{\bfseries}]
    \item[3.1.14] Система должна обеспечивать полнотекстовый поиск по всем разделам портала.
    \item[3.1.15] Система должна поддерживать автодополнение при вводе запроса.
    \item[3.1.16] Система должна отображать подсказки на основе популярных запросов.
    \item[3.1.17] Система должна позволять фильтрацию результатов по типу контента (праздник, имя, событие, гороскоп, новость, статья).
    \item[3.1.18] Система должна сохранять историю поиска для авторизованных пользователей.
    \item[3.1.19] Система должна позволять регистрацию пользователей через электронную почту.
    \item[3.1.20] Система должна позволять авторизацию по логину/паролю или через социальные сети.
    \item[3.1.21] Система должна предоставлять функцию восстановления пароля через электронную почту.
    \item[3.1.22] Система должна позволять добавление контента в избранное.
    \item[3.1.23] Система должна позволять настройку персонального гороскопа (выбор знака зодиака).
    \item[3.1.24] Система должна предоставлять возможность настройки уведомлений о праздниках и именинах.
\end{description}

\subsubsection*{Редакционный контент}
\begin{description}[leftmargin=!,labelwidth=\widthof{\bfseries}]
    \item[3.1.25] Система должна позволять редактору создавать публикации (новости, статьи, анонсы) в статусе «черновик».
    \item[3.1.26] Система должна поддерживать жизненный цикл публикации: черновик, на модерации, опубликовано, архив.
    \item[3.1.27] Система должна поддерживать отложенную публикацию по заданным дате и времени.
    \item[3.1.28] Система должна вести историю изменений публикаций с указанием автора правки и времени изменения.
    \item[3.1.29] Система должна предоставлять возможность отката публикации к одной из предыдущих версий.
    \item[3.1.30] Система должна позволять редактору управлять SEO-параметрами публикации (title, description, slug, keywords).
\end{description}

\subsubsection*{Реклама и монетизация}
\begin{description}[leftmargin=!,labelwidth=\widthof{\bfseries}]
    \item[3.1.31] Система должна поддерживать размещение рекламных блоков в заданных областях страниц (header, sidebar, inline, footer).
    \item[3.1.32] Система должна позволять настраивать период активности рекламной кампании (дата начала и дата завершения).
    \item[3.1.33] Система должна поддерживать таргетинг рекламных блоков по разделам сайта и типам устройств.
    \item[3.1.34] Система должна собирать статистику показов, кликов и CTR по рекламным кампаниям.
    \item[3.1.35] Система должна маркировать рекламные материалы и учитывать пользовательские согласия для персонализированной рекламы.
\end{description}

\section{Nonfunctional Requirements}

\subsection{Usability Requirements}
\begin{description}[leftmargin=!,labelwidth=\widthof{\bfseries}]
    \item[4.1.1] Система должна корректно отображаться на устройствах с разрешением экрана от 320px (мобильные) до 3840px (4K).
    \item[4.1.2] Система должна поддерживать масштабирование интерфейса от 50\% до 200\% без потери функциональности.
    \item[4.1.3] Текстовый контент и элементы управления должны оставаться читаемыми и доступными при масштабе 200\%.
    \item[4.1.4] Система должна обеспечивать навигацию с клавиатуры для пользователей с ограниченными возможностями.
\end{description}

\subsection{Reliability Requirements}
\begin{description}[leftmargin=!,labelwidth=\widthof{\bfseries}]
    \item[4.2.1] Доступность системы должна составлять не менее 99.5\% в течение календарного месяца.
    \item[4.2.2] Среднее время восстановления после сбоя не должно превышать 30 минут.
    \item[4.2.3] Система должна выполнять ежедневное резервное копирование базы данных.
    \item[4.2.4] Резервные копии должны храниться не менее 30 дней.
    \item[4.2.5] Система должна уведомлять администратора о критических сбоях в течение 5 минут.
\end{description}

\subsection{Performance Requirements}
\begin{description}[leftmargin=!,labelwidth=\widthof{\bfseries}]
    \item[4.3.1] Время загрузки главной страницы не должно превышать 2 секунды при скорости соединения 10 Мбит/с.
    \item[4.3.2] Время ответа сервера на запрос календаря не должно превышать 300 мс.
    \item[4.3.3] Система должна поддерживать до 5000 одновременных пользователей без деградации производительности.
    \item[4.3.4] Время выполнения поискового запроса не должно превышать 1 секунды.
\end{description}

\subsection{Licensing Requirements}
\begin{description}[leftmargin=!,labelwidth=\widthof{\bfseries}]
    \item[4.4.1] Исходный код системы должен распространяться под закрытой коммерческой лицензией.
    \item[4.4.2] Использование сторонних библиотек допускается только при наличии совместимых лицензий.
\end{description}
